\chapter{Edward Witten --- Fields Medal 1990}\label{ch1990edwardwitten:edward-witten-fields-medal-1990-a-comprehensive-report}

\begin{center}\rule{0.5\linewidth}{0.5pt}\end{center}

\textbf{Prepared:} July 2026 \textbf{Subject:} Edward Witten (b. August
26, 1951, Baltimore, Maryland, USA) \textbf{Position:} Professor
Emeritus, Institute for Advanced Study, Princeton (since 1987)
\textbf{Doctoral Advisor:} David Gross (Princeton University)
\textbf{Fields Medal Citation:} \emph{``For his contributions to
mathematical physics, including the proof of the positive energy
theorem, supersymmetry, and the Jones polynomial in quantum field
theory.''}

\begin{center}\rule{0.5\linewidth}{0.5pt}\end{center}

\section{1. Introduction}\label{ch1990edwardwitten:introduction}

Edward Witten was awarded the Fields Medal in 1990 --- the first and
only physicist to receive the Fields Medal. His work has substantially
reshaped both theoretical physics and mathematics, building bridges
between the two fields that have transformed each.

Witten's contributions include: - \textbf{The positive energy theorem}
(with Richard Schoen and Shing-Tung Yau). - \textbf{Supersymmetry and
supergravity.} - \textbf{The Jones polynomial in Chern--Simons theory.}
- \textbf{String theory and M-theory.} - \textbf{Mirror symmetry and
Calabi--Yau compactifications.} - \textbf{Seiberg--Witten invariants.} -
\textbf{Topological quantum field theory.}

\begin{center}\rule{0.5\linewidth}{0.5pt}\end{center}

\section{2. Biographical Background}\label{ch1990edwardwitten:biographical-background}

\subsection{2.1 Early Life}\label{ch1990edwardwitten:early-life}

Edward Witten was born in Baltimore in 1951. His father, Louis Witten,
was a theoretical physicist specializing in general relativity. Edward
studied at Brandeis University and the University of Wisconsin--Madison.

\subsection{2.2 Education}\label{ch1990edwardwitten:education}

\begin{itemize}
\tightlist
\item
 \textbf{1971:} BA in History, Brandeis University.
\item
 \textbf{1974:} MA in Mathematics, University of Wisconsin--Madison.
\item
 \textbf{1976:} PhD in Physics, Princeton University (under David
 Gross).
\end{itemize}

\subsection{2.3 Academic Career}\label{ch1990edwardwitten:academic-career}

\begin{itemize}
\tightlist
\item
 \textbf{1976--1977:} Postdoctoral Fellow, Harvard University.
\item
 \textbf{1977--1980:} Junior Fellow, Harvard Society of Fellows.
\item
 \textbf{1980--1987:} Professor of Physics, Princeton University.
\item
 \textbf{1987--present:} Professor, Institute for Advanced Study,
 Princeton.
\item
 \textbf{1990:} Fields Medal.
\item
 \textbf{1991:} MacArthur Fellowship.
\item
 \textbf{2004:} Appointed Charles Simonyi Professor, IAS.
\item
 \textbf{2008:} Awarded the Crafoord Prize.
\end{itemize}

\begin{center}\rule{0.5\linewidth}{0.5pt}\end{center}

\section{3. The Positive Energy
Theorem}\label{ch1990edwardwitten:the-positive-energy-theorem}

\subsection{3.1 Background}\label{ch1990edwardwitten:background}

In general relativity, the \textbf{positive energy theorem} states that
the total energy (mass) of an asymptotically flat spacetime is
non-negative, and zero only for flat Minkowski space.

\subsection{3.2 Witten's Proof}\label{ch1990edwardwitten:wittens-proof}

Witten (with Schoen and Yau) gave a proof using spinors: 1. Construct a
spinor field satisfying the Dirac equation. 2. Express the total mass as
an integral over the boundary. 3. The integral is non-negative due to
the dominant energy condition.

\subsection{3.3 Significance}\label{ch1990edwardwitten:significance}

This theorem is fundamental to general relativity, ensuring that
gravitationally bound systems have positive energy.

\begin{center}\rule{0.5\linewidth}{0.5pt}\end{center}

\section{4. Supersymmetry}\label{ch1990edwardwitten:supersymmetry}

\subsection{4.1 Definition}\label{ch1990edwardwitten:definition}

\textbf{Supersymmetry} is a symmetry relating bosons (integer spin) and
fermions (half-integer spin). It is a central concept in modern
theoretical physics.

\subsection{4.2 Witten's Contributions}\label{ch1990edwardwitten:wittens-contributions}

Witten has made numerous contributions to supersymmetry: -
\textbf{Supersymmetric quantum mechanics} as a toy model. -
\textbf{Supersymmetric gauge theories} in various dimensions. -
\textbf{The index theorem and supersymmetry.} - \textbf{Dynamical
supersymmetry breaking.}

\subsection{4.3 Supersymmetric Gauge
Theories}\label{ch1990edwardwitten:supersymmetric-gauge-theories}

Witten studied the dynamics of supersymmetric gauge theories: - The
moduli space of vacua. - The chiral ring and its properties. - The
relationship between supersymmetry and confinement.

\begin{center}\rule{0.5\linewidth}{0.5pt}\end{center}

\section{5. The Jones Polynomial and Chern--Simons
Theory}\label{ch1990edwardwitten:the-jones-polynomial-and-chernsimons-theory}

\subsection{5.1 Chern--Simons Theory}\label{ch1990edwardwitten:chernsimons-theory}

Witten showed that the \textbf{Jones polynomial} arises naturally from
\textbf{Chern--Simons gauge theory} in three dimensions.

\subsection{5.2 The Partition Function}\label{ch1990edwardwitten:the-partition-function}

The partition function of Chern--Simons theory with gauge group
\(SU(2)\) at level \(k\) is:

\[
Z(M) = \int \mathcal{D}A \, \exp\left(\frac{ik}{4\pi} \int_M \text{Tr}\left(A \wedge dA + \frac{2}{3} A \wedge A \wedge A\right)\right)
\]

\subsection{5.3 The Wilson Loop
Expectation}\label{ch1990edwardwitten:the-wilson-loop-expectation}

Witten showed that the expectation value of a Wilson loop along a knot
\(K\):

\[
\langle W_R(K) \rangle = \frac{\int \mathcal{D}A \, \text{Tr}_R(P e^{\oint_K A}) \, e^{iS_{CS}(A)}}{\int \mathcal{D}A \, e^{iS_{CS}(A)}}
\]

is equal to the Jones polynomial (or its generalization, the HOMFLY
polynomial) evaluated at \(q = e^{2\pi i/(k+2)}\).

\subsection{5.4 Three-Manifold
Invariants}\label{ch1990edwardwitten:three-manifold-invariants}

Witten constructed invariants of three-manifolds from Chern--Simons
theory, now called \textbf{Witten--Reshetikhin--Turaev invariants}.

\subsection{5.5 Impact}\label{ch1990edwardwitten:impact}

This work connected: - Knot theory and quantum field theory. - The Jones
polynomial and conformal field theory. - Three-manifold topology and
quantum groups.

\begin{center}\rule{0.5\linewidth}{0.5pt}\end{center}

\section{6. Topological Quantum Field
Theory}\label{ch1990edwardwitten:topological-quantum-field-theory}

\subsection{6.1 Definition}\label{ch1990edwardwitten:definition-1}

Witten introduced \textbf{topological quantum field theory} (TQFT) as a
quantum field theory whose correlation functions are topological
invariants.

\subsection{6.2 Examples}\label{ch1990edwardwitten:examples}

\begin{itemize}
\tightlist
\item
 \textbf{Chern--Simons theory:} A three-dimensional TQFT.
\item
 \textbf{Topological sigma models:} Two-dimensional TQFTs related to
 mirror symmetry.
\item
 \textbf{Donaldson theory:} A four-dimensional TQFT related to
 Donaldson invariants.
\end{itemize}

\subsection{6.3 The Atiyah--Segal Axioms}\label{ch1990edwardwitten:the-atiyahsegal-axioms}

Witten's work inspired the axiomatization of TQFT by Atiyah and Segal.

\begin{center}\rule{0.5\linewidth}{0.5pt}\end{center}

\section{7. Seiberg--Witten Theory}\label{ch1990edwardwitten:seibergwitten-theory}

\subsection{7.1 The Seiberg--Witten
Equations}\label{ch1990edwardwitten:the-seibergwitten-equations}

In 1994, Witten (with Nathan Seiberg) introduced the
\textbf{Seiberg--Witten equations}, a new set of equations in
four-dimensional gauge theory:

\[
\not{D}_A \psi = 0, \quad F_A^+ = \psi \otimes \psi^* - \frac{1}{2}|\psi|^2 I
\]

\subsection{7.2 Seiberg--Witten
Invariants}\label{ch1990edwardwitten:seibergwitten-invariants}

The Seiberg--Witten equations give invariants of smooth four-manifolds
that are much easier to compute than Donaldson invariants.

\subsection{7.3 Applications}\label{ch1990edwardwitten:applications}

\begin{itemize}
\tightlist
\item
 \textbf{Classification of four-manifolds:} The Seiberg--Witten
 invariants distinguish many smooth structures.
\item
 \textbf{The Thom conjecture:} Proved using Seiberg--Witten theory.
\item
 \textbf{Symplectic geometry:} Applications to the topology of
 symplectic four-manifolds.
\end{itemize}

\begin{center}\rule{0.5\linewidth}{0.5pt}\end{center}

\section{8. String Theory}\label{ch1990edwardwitten:string-theory}

\subsection{8.1 Background}\label{ch1990edwardwitten:background-1}

\textbf{String theory} posits that fundamental particles are
one-dimensional ``strings'' rather than point particles.

\subsection{8.2 Witten's Contributions}\label{ch1990edwardwitten:wittens-contributions-1}

Witten has been a leading figure in string theory: - \textbf{Heterotic
string theory:} The construction of the heterotic string. -
\textbf{Calabi--Yau compactifications:} The study of string theory
compactified on Calabi--Yau manifolds. - \textbf{Mirror symmetry:} The
duality between Calabi--Yau manifolds and their mirrors. -
\textbf{D-branes:} The discovery of D-branes and their role in string
dualities.

\begin{center}\rule{0.5\linewidth}{0.5pt}\end{center}

\section{9. M-Theory}\label{ch1990edwardwitten:m-theory}

\subsection{9.1 The Unification of String
Theories}\label{ch1990edwardwitten:the-unification-of-string-theories}

In 1995, Witten proposed \textbf{M-theory} as the unifying framework for
the five consistent superstring theories: - Type I - Type IIA - Type IIB
- Heterotic \(SO(32)\) - Heterotic \(E_8 \times E_8\)

\subsection{9.2 Duality}\label{ch1990edwardwitten:duality}

M-theory lives in 11 dimensions and is related to 10-dimensional string
theories through dualities: - \textbf{T-duality:} Equivalence of string
theories on large and small circles. - \textbf{S-duality:} Equivalence
of strongly and weakly coupled theories. - \textbf{U-duality:} The
unification of T and S dualities.

\begin{center}\rule{0.5\linewidth}{0.5pt}\end{center}

\section{10. Mirror Symmetry}\label{ch1990edwardwitten:mirror-symmetry}

\subsection{10.1 The Conjecture}\label{ch1990edwardwitten:the-conjecture}

\textbf{Mirror symmetry} predicts that for every Calabi--Yau threefold
\(X\), there exists a mirror Calabi--Yau threefold \(\hat{X}\) such that

\[
h^{1,1}(X) = h^{2,1}(\hat{X}), \quad h^{2,1}(X) = h^{1,1}(\hat{X})
\]

\subsection{10.2 Witten's Contributions}\label{ch1990edwardwitten:wittens-contributions-2}

Witten contributed: - The connection between mirror symmetry and
T-duality. - The sigma model approach to mirror symmetry. - The quantum
cohomology ring.

\begin{center}\rule{0.5\linewidth}{0.5pt}\end{center}

\section{11. The Index Theorem and
Anomalies}\label{ch1990edwardwitten:the-index-theorem-and-anomalies}

\subsection{11.1 The Atiyah--Singer Index
Theorem}\label{ch1990edwardwitten:the-atiyahsinger-index-theorem}

The \textbf{Atiyah--Singer index theorem} relates the analytical index
of an elliptic operator to a topological invariant.

\subsection{11.2 Witten's Contributions}\label{ch1990edwardwitten:wittens-contributions-3}

Witten contributed to: - The proof of the index theorem using
supersymmetry. - The study of gravitational anomalies. - The
relationship between anomalies and the topology of spacetime.

\begin{center}\rule{0.5\linewidth}{0.5pt}\end{center}

\section{12. The Work on Gravitational
Anomalies}\label{ch1990edwardwitten:the-work-on-gravitational-anomalies}

\subsection{12.1 Definition}\label{ch1990edwardwitten:definition-2}

\textbf{Gravitational anomalies} arise when a quantum field theory that
is classically invariant under diffeomorphisms fails to be invariant
after quantization.

\subsection{12.2 Witten's Classification}\label{ch1990edwardwitten:wittens-classification}

Witten classified gravitational anomalies and showed that they obstruct
the consistency of quantum gravity theories.

\begin{center}\rule{0.5\linewidth}{0.5pt}\end{center}

\section{13. The Work on the BPS Bound}\label{ch1990edwardwitten:the-work-on-the-bps-bound}

\subsection{13.1 Definition}\label{ch1990edwardwitten:definition-3}

The \textbf{Bogomolny--Prasad--Sommerfield (BPS) bound} is a lower bound
on the mass of solitons in supersymmetric theories.

\subsection{13.2 Witten's Contributions}\label{ch1990edwardwitten:wittens-contributions-4}

Witten studied: - BPS states in supersymmetric gauge theories. - The
relationship between BPS states and moduli spaces. - The wall-crossing
phenomenon for BPS states.

\begin{center}\rule{0.5\linewidth}{0.5pt}\end{center}

\section{14. The Work on the Donaldson
Invariants}\label{ch1990edwardwitten:the-work-on-the-donaldson-invariants}

\subsection{14.1 Donaldson Theory}\label{ch1990edwardwitten:donaldson-theory}

\textbf{Donaldson invariants} are invariants of smooth four-manifolds
obtained from the moduli space of anti-self-dual connections.

\subsection{14.2 Witten's Contributions}\label{ch1990edwardwitten:wittens-contributions-5}

Witten reformulated Donaldson theory as a topological quantum field
theory: - The Donaldson--Witten theory. - The relationship between
Donaldson and Seiberg--Witten invariants. - The computation of Donaldson
invariants using Seiberg--Witten theory.

\begin{center}\rule{0.5\linewidth}{0.5pt}\end{center}

\section{15. The Work on the Geometric Langlands
Program}\label{ch1990edwardwitten:the-work-on-the-geometric-langlands-program}

\subsection{15.1 Background}\label{ch1990edwardwitten:background-2}

The \textbf{geometric Langlands program} is a geometric reformulation of
the Langlands conjectures in number theory.

\subsection{15.2 Witten's Contributions}\label{ch1990edwardwitten:wittens-contributions-6}

Witten (with Kapustin) showed: - The geometric Langlands correspondence
arises from S-duality in \(N=4\) super Yang--Mills theory. - The Hecke
eigensheaves correspond to boundary conditions in the quantum field
theory. - This provides a physical origin for the Langlands program.

\begin{center}\rule{0.5\linewidth}{0.5pt}\end{center}

\section{16. The Work on Khovanov
Homology}\label{ch1990edwardwitten:the-work-on-khovanov-homology}

\subsection{16.1 Definition}\label{ch1990edwardwitten:definition-4}

\textbf{Khovanov homology} is a categorification of the Jones
polynomial.

\subsection{16.2 Witten's Contributions}\label{ch1990edwardwitten:wittens-contributions-7}

Witten (with Gukov and others) studied: - The physical interpretation of
Khovanov homology in string theory. - The relationship between Khovanov
homology and BPS states. - The categorification of knot invariants from
M-theory.

\begin{center}\rule{0.5\linewidth}{0.5pt}\end{center}

\section{17. The Work on the Space of
Vacua}\label{ch1990edwardwitten:the-work-on-the-space-of-vacua}

\subsection{17.1 Moduli Spaces}\label{ch1990edwardwitten:moduli-spaces}

Witten studied the moduli space of vacua in supersymmetric gauge
theories: - The classical moduli space is parameterized by expectation
values of scalar fields. - Quantum effects modify the moduli space and
can break supersymmetry.

\begin{center}\rule{0.5\linewidth}{0.5pt}\end{center}

\section{18. The Work on the AdS/CFT
Correspondence}\label{ch1990edwardwitten:the-work-on-the-adscft-correspondence}

\subsection{18.1 The Correspondence}\label{ch1990edwardwitten:the-correspondence}

The \textbf{AdS/CFT correspondence} (Maldacena, 1997) relates string
theory in anti-de Sitter space to conformal field theory on the
boundary.

\subsection{18.2 Witten's Contributions}\label{ch1990edwardwitten:wittens-contributions-8}

Witten significantly advanced the AdS/CFT correspondence: - Holographic
entanglement entropy. - The computation of correlation functions in
AdS/CFT. - The relationship between AdS/CFT and the renormalization
group.

\begin{center}\rule{0.5\linewidth}{0.5pt}\end{center}

\section{19. The Work on the Verlinde
Formula}\label{ch1990edwardwitten:the-work-on-the-verlinde-formula}

\subsection{19.1 The Formula}\label{ch1990edwardwitten:the-formula}

The \textbf{Verlinde formula} computes the dimensions of the spaces of
conformal blocks in rational conformal field theory.

\subsection{19.2 Witten's Contributions}\label{ch1990edwardwitten:wittens-contributions-9}

Witten gave a derivation of the Verlinde formula from Chern--Simons
theory: - The Hilbert space of Chern--Simons theory is the space of
conformal blocks. - The Verlinde formula computes the dimension of this
Hilbert space.

\begin{center}\rule{0.5\linewidth}{0.5pt}\end{center}

\section{20. The Work on the Analytic Continuation in Chern--Simons
Theory}\label{ch1990edwardwitten:the-work-on-the-analytic-continuation-in-chernsimons-theory}

\subsection{20.1 The Problem}\label{ch1990edwardwitten:the-problem}

Chern--Simons theory is defined for integer levels, but can be
analytically continued to complex levels.

\subsection{20.2 Witten's Contributions}\label{ch1990edwardwitten:wittens-contributions-10}

Witten studied: - The analytic continuation of Chern--Simons theory. -
The relationship to the volume conjecture. - The asymptotic expansion of
the Jones polynomial.

\begin{center}\rule{0.5\linewidth}{0.5pt}\end{center}

\section{21. The Work on the Volume
Conjecture}\label{ch1990edwardwitten:the-work-on-the-volume-conjecture}

\subsection{21.1 The Conjecture}\label{ch1990edwardwitten:the-conjecture-1}

The \textbf{volume conjecture} relates the asymptotic behavior of the
colored Jones polynomial of a hyperbolic knot to the hyperbolic volume
of the knot complement.

\subsection{21.2 Witten's Contributions}\label{ch1990edwardwitten:wittens-contributions-11}

Witten contributed to the physical understanding of the volume
conjecture: - The relationship between Chern--Simons theory and
hyperbolic geometry. - The asymptotic expansion of the path integral.

\begin{center}\rule{0.5\linewidth}{0.5pt}\end{center}

\section{22. The Work on the Three-Manifold
Invariants}\label{ch1990edwardwitten:the-work-on-the-three-manifold-invariants}

\subsection{22.1 Witten--Reshetikhin--Turaev
Invariants}\label{ch1990edwardwitten:wittenreshetikhinturaev-invariants}

Witten constructed invariants of three-manifolds from Chern--Simons
theory. These were later put on a rigorous mathematical foundation by
Reshetikhin and Turaev.

\begin{center}\rule{0.5\linewidth}{0.5pt}\end{center}

\section{23. The Work on Quantum
Cohomology}\label{ch1990edwardwitten:the-work-on-quantum-cohomology}

\subsection{23.1 Definition}\label{ch1990edwardwitten:definition-5}

\textbf{Quantum cohomology} is a deformation of the ordinary cohomology
ring of a symplectic manifold, using Gromov--Witten invariants.

\subsection{23.2 Witten's Contributions}\label{ch1990edwardwitten:wittens-contributions-12}

Witten contributed to: - The relationship between quantum cohomology and
mirror symmetry. - The sigma model approach to quantum cohomology. - The
computation of Gromov--Witten invariants.

\begin{center}\rule{0.5\linewidth}{0.5pt}\end{center}

\section{24. The Work on Gromov--Witten
Invariants}\label{ch1990edwardwitten:the-work-on-gromovwitten-invariants}

\subsection{24.1 Definition}\label{ch1990edwardwitten:definition-6}

\textbf{Gromov--Witten invariants} count curves on a symplectic
manifold.

\subsection{24.2 Witten's Contributions}\label{ch1990edwardwitten:wittens-contributions-13}

Witten studied: - The relationship between Gromov--Witten invariants and
quantum cohomology. - The topological sigma model and Gromov--Witten
theory. - The application of Gromov--Witten invariants to enumerative
geometry.

\begin{center}\rule{0.5\linewidth}{0.5pt}\end{center}

\section{25. Awards and Recognition}\label{ch1990edwardwitten:awards-and-recognition}

\textbf{1990:} - Fields Medal (International Mathematical Union)

\textbf{1991:} - MacArthur Fellowship

\textbf{2004:} - Henri Poincaré Prize

\textbf{2008:} - Crafoord Prize in Astronomy

\textbf{2010:} - Nemmers Prize in Mathematics

\textbf{2012:} - Fundamental Physics Prize

\textbf{2014:} - Lorentz Medal

\textbf{2016:} - Breakthrough Prize in Fundamental Physics

\textbf{2017:} - Election to the National Academy of Sciences

\begin{center}\rule{0.5\linewidth}{0.5pt}\end{center}

\section{26. Selected Major
Publications}\label{ch1990edwardwitten:selected-major-publications}

\begin{enumerate}
\def\labelenumi{\arabic{enumi}.}
\item
 \textbf{``A new proof of the positive energy theorem.''}
 \emph{Communications in Mathematical Physics} 80, 381--402, 1981. ---
 \emph{The proof of the positive energy theorem.}
\item
 \textbf{``Supersymmetry and the index theorem.''} \emph{Journal of
 Differential Geometry} 17, 661--692, 1982. --- \emph{The index theorem
 from supersymmetry.}
\item
 \textbf{``Quantum field theory and the Jones polynomial.''}
 \emph{Communications in Mathematical Physics} 121, 351--399, 1989. ---
 \emph{Chern--Simons theory and the Jones polynomial.}
\item
 \textbf{``Topological quantum field theory.''} \emph{Communications in
 Mathematical Physics} 117, 353--386, 1988. --- \emph{The introduction
 of TQFT.}
\item
 \textbf{``Monopoles and four-manifolds.''} (with N. Seiberg)
 \emph{Mathematical Research Letters} 1, 769--796, 1994. ---
 \emph{Seiberg--Witten theory.}
\item
 \textbf{``String theory dynamics in various dimensions.''}
 \emph{Nuclear Physics B} 443, 85--126, 1995. --- \emph{The proposal of
 M-theory.}
\item
 \textbf{``Mirror symmetry and the moduli space of Calabi--Yau
 manifolds.''} \emph{Nuclear Physics B} 403, 1993. --- \emph{Mirror
 symmetry and string theory.}
\end{enumerate}

\begin{center}\rule{0.5\linewidth}{0.5pt}\end{center}

\section{27. References and Further
Reading}\label{ch1990edwardwitten:references-and-further-reading}

\subsection{Official Sources}\label{ch1990edwardwitten:official-sources}

\begin{itemize}
\tightlist
\item
 IMU Fields Medal 1990 Announcement:
 https://www.mathunion.org/imu-awards/fields-medal/fields-medal-1990
\end{itemize}

\subsection{Key Technical References}\label{ch1990edwardwitten:key-technical-references}

\begin{itemize}
\tightlist
\item
 Witten, E. ``Quantum field theory and the Jones polynomial.''
 \emph{Commun. Math. Phys.} 121, 1989.
\item
 Witten, E. ``Topological quantum field theory.'' \emph{Commun. Math.
 Phys.} 117, 1988.
\item
 Seiberg, N. and Witten, E. ``Monopoles and four-manifolds.''
 \emph{Math. Res. Lett.} 1, 1994.
\end{itemize}

\subsection{Textbooks}\label{ch1990edwardwitten:textbooks}

\begin{itemize}
\tightlist
\item
 Deligne, P. et al.~\emph{Quantum Fields and Strings: A Course for
 Mathematicians}. AMS, 1999.
\item
 Freed, D. and Uhlenbeck, K. \emph{Instantons and Four-Manifolds}.
 Springer, 1984.
\item
 Kaku, M. \emph{Strings, Conformal Fields, and M-Theory}. Springer,
 2000.
\end{itemize}

\subsection{Biography}\label{ch1990edwardwitten:biography}

\begin{itemize}
\tightlist
\item
 Atiyah, M. ``The Work of Edward Witten.'' ICM 1990 Laudatio.
\item
 Greene, B. \emph{The Elegant Universe}. Norton, 1999.
\end{itemize}

\begin{center}\rule{0.5\linewidth}{0.5pt}\end{center}

\emph{(End of report)}
